% Zone „Das kennst du schon“ – aus bank/quadratische-gleichungen/zone.jsonl (zusammenbau v0.1)
\einheitenkopf[zone][Kennst du schon]{Das kennst du schon}
\setcounter{aufgabe}{0}

%% TODO Titel: Ich-kann-Satz fehlt in der Bank (Platzhalter: Kettenname) – Nr. 1
%% TODO Anweisung: ein Satz über der ersten Teilaufgabe fehlt in der Bank – Nr. 1
\begin{aufgabe}{Quadrieren auch negativer Zahlen und Dezimalzahlen, Quadrat vor Punkt vor Strich}
\begin{teile}
\teil $7^2$ – wie viel? \leerfeld
\teil $(-1{,}2)^2$ – wie viel? \leerfeld
\end{teile}
\end{aufgabe}

%% TODO Titel: Ich-kann-Satz fehlt in der Bank (Platzhalter: Kettenname) – Nr. 2
%% TODO Anweisung: ein Satz über der ersten Teilaufgabe fehlt in der Bank – Nr. 2
\begin{aufgabe}{Quadratwurzel ziehen}
\begin{teile}
\teil $\sqrt{49}$ – wie viel? \leerfeld
\teil $\sqrt{11}$ – auf zwei Stellen? \leerfeld
\end{teile}
\end{aufgabe}

%% TODO Titel: Ich-kann-Satz fehlt in der Bank (Platzhalter: Kettenname) – Nr. 3
%% TODO Anweisung: ein Satz über der ersten Teilaufgabe fehlt in der Bank – Nr. 3
\begin{aufgabe}{Lineare Gleichung mit Äquivalenzumformungen und Strich lösen, auch negative Vorzahl und x auf beiden Seiten}
\begin{gleichungsraster}[2]
\gl{\text{x + 8 = 15}} &
\gl{\text{-2x = 14}} \\
\end{gleichungsraster}
\end{aufgabe}

%% TODO Titel: Ich-kann-Satz fehlt in der Bank (Platzhalter: Kettenname) – Nr. 4
%% TODO Anweisung: ein Satz über der ersten Teilaufgabe fehlt in der Bank – Nr. 4
\begin{aufgabe}{Lösung durch Einsetzen prüfen mit (wA)/(fA), auch mit negativer Zahl}
\begin{teile}
\teil Ist $x = 4$ Lösung von $x + 9 = 13$? \janein
\teil Ist $x = -6$ Lösung von $x + 10 = 4$? \janein
\end{teile}
\end{aufgabe}

%% TODO Titel: Ich-kann-Satz fehlt in der Bank (Platzhalter: Kettenname) – Nr. 5
%% TODO Anweisung: ein Satz über der ersten Teilaufgabe fehlt in der Bank – Nr. 5
\begin{aufgabe}{Klammer zuerst und Punkt vor Strich mit negativen Zahlen}
\begin{teile}
\teil $3 \cdot (2 + 4)$ – wie viel? \leerfeld
\teil $(-4) \cdot (-4 + 9)$ – wie viel? \leerfeld
\end{teile}
\end{aufgabe}

%% TODO Titel: Ich-kann-Satz fehlt in der Bank (Platzhalter: Kettenname) – Nr. 6
%% TODO Anweisung: ein Satz über der ersten Teilaufgabe fehlt in der Bank – Nr. 6
\begin{aufgabe}{Scheitelpunktform lesen}
\begin{teile}
\teil $y = (x - 2)^2 + 1$ – Scheitel? S( \leerfeld | \leerfeld )
\teil $y = (x + 1)^2 - 6$ – Scheitel? S( \leerfeld | \leerfeld )
\end{teile}
\end{aufgabe}

%% TODO Titel: Ich-kann-Satz fehlt in der Bank (Platzhalter: Kettenname) – Nr. 7
%% TODO Anweisung: ein Satz über der ersten Teilaufgabe fehlt in der Bank – Nr. 7
\begin{aufgabe}{Ausklammern, Ausmultiplizieren und binomische Formeln}
\begin{teile}
\teil $x \cdot (x + 4)$ – ausmultiplizieren? \leerfeld
\teil $(x - 6)^2$ – ausmultiplizieren? \leerfeld
\end{teile}
\end{aufgabe}

%% TODO Titel: Ich-kann-Satz fehlt in der Bank (Platzhalter: Kettenname) – Nr. 8
%% TODO Anweisung: ein Satz über der ersten Teilaufgabe fehlt in der Bank – Nr. 8
\begin{aufgabe}{Terme ordnen und zusammenfassen, x² zuerst, dann x-Glieder, dann Zahlen}
\begin{teile}
\teil $3x + x^2 + 2$ – ordnen? \leerfeld
\teil $x^2 - 8 + 3x + 2$ – ordnen und zusammenfassen? \leerfeld
\end{teile}
\end{aufgabe}

%% TODO Titel: Ich-kann-Satz fehlt in der Bank (Platzhalter: Kettenname) – Nr. 9
%% TODO Anweisung: ein Satz über der ersten Teilaufgabe fehlt in der Bank – Nr. 9
\begin{aufgabe}{Quadrieren auch negativer Zahlen und Dezimalzahlen, Quadrat vor Punkt vor Strich – Fehler finden}
\begin{teile}
\teil Wo ist der Fehler? $(-8)^2 = -64$
\end{teile}
\end{aufgabe}

%% TODO Titel: Ich-kann-Satz fehlt in der Bank (Platzhalter: Kettenname) – Nr. 10
%% TODO Anweisung: ein Satz über der ersten Teilaufgabe fehlt in der Bank – Nr. 10
\begin{aufgabe}{Quadrieren auch negativer Zahlen und Dezimalzahlen, Quadrat vor Punkt vor Strich – selbst rechnen}
\begin{teile}
\teil $(-13)^2$ – wie viel? \leerfeld
\end{teile}
\end{aufgabe}
